\chapter{Frameworks for Service-Oriented Applications}
\label{chap:soa}

\section{Purpose of This Review}

This chapter is not intended to be a broad survey of every service framework.
Instead, it reviews what existing systems have already achieved and identifies
the open issues that motivate NDNSF. The key question is not whether systems
such as gRPC, ROS, microservices, or NDN Service Call are useful; they clearly
are. The key question is why dynamic, secure, collaborative applications still
require another framework despite the existence of those systems.

\section{Host-Centric RPC and Microservice Frameworks}

RPC systems such as gRPC~\cite{grpc}, HTTP/REST systems~\cite{fielding2000rest},
and microservice platforms have become the dominant way to build distributed
software. They provide typed APIs, transport bindings, load balancing, service
registration, and production tooling. These systems work very well when
services can be identified by stable endpoints, when channels can be
established and maintained, and when security can be expressed as channel
security plus service-side authorization.

However, the abstraction becomes less natural for applications with mobility
and intermittent connectivity. A user of a UAV video service may care about the
latest frames from drone A, not whether drone A is currently reachable through
a particular IP address. A distributed inference client may care about the
service name for a model stage, not a fixed GPU server. A mission-planning
application may need any available drone that satisfies a role, not exactly one
preconfigured endpoint. Host-centric frameworks can support these scenarios,
but the burden shifts to external discovery, retry, state tracking, and
application-specific coordination.

\subsection{What These Systems Solve}

RPC and microservice systems solve several problems that NDNSF should not
ignore. First, they give developers a familiar programming model: a service has
an interface, a request type, a response type, error handling, and often a
deadline. Second, production systems around gRPC and REST provide mature
tooling for load balancing, logging, monitoring, deployment, and versioning.
Third, channel security such as TLS gives a practical and widely deployed way
to protect client-server communication. For stable cloud services, these
features are often sufficient.

NDNSF does not attempt to replace this engineering knowledge. Instead, it asks
which assumptions break when service providers are mobile, intermittently
connected, replicated, or embedded in resource-constrained edge environments.
In these settings, a service call is not only a channel operation. It is also a
decision about which provider is currently authorized, reachable, capable, and
worth selecting. Existing host-centric frameworks can implement this behavior,
but the logic is typically distributed across service discovery systems,
load-balancers, application retries, middleware plugins, and service-specific
policy code.

\subsection{Open Issues for NDNSF}

The first issue is endpoint dependence. Even if modern service meshes hide IP
addresses behind logical service names, the underlying communication still
depends on reaching an endpoint through host-centric routing. The second issue
is data lifetime. Once a response, video chunk, model shard, or telemetry log is
stored outside the original channel, applications must add their own
verification and encryption model. The third issue is provider multiplicity.
RPC normally selects a single server before the request is sent; NDNSF instead
allows multiple providers to observe a request, decide whether to acknowledge
it, and expose metadata that helps the user choose. These differences motivate
a data-centric service framework rather than another IP endpoint wrapper.

\section{Robotics Middleware and Service Composition}

Robotics middleware, including ROS and ROS2~\cite{ros,ros2}, provides a more
application-aware model for robot components. Nodes can publish topics, call
services, and interact through standardized message types. ROS2 improves
quality-of-service control and deployment flexibility compared with earlier
ROS systems. These frameworks are important precedents for UAV applications
because they treat sensors, control, telemetry, and perception as composable
modules.

The remaining gap is that robotics middleware is not a general data-centric
network security framework. Trust anchors, certificate deployment, data-at-rest
protection, multi-provider service selection, and cross-domain service
authorization are not the main abstractions. In addition, a UAV application
deployed across a ground station, multiple drones, and intermittent wireless
links still needs naming, verification, and recovery mechanisms that are robust
outside a single local robotics domain.

\subsection{Lessons for UAV Service Containers}

ROS is particularly important for this dissertation because it shows that
robotics applications are naturally service-oriented. A UAV system is not a
single communication flow. It contains flight-control commands, telemetry
streams, camera frames, mapping, mission planning, object detection, logging,
and health monitoring. ROS-style composition validates the idea that UAV
software should be decomposed into services.

However, NDNSF targets a different layer. ROS and ROS2 are strong middleware
systems for local and distributed robot components, while NDNSF focuses on
named, secure, network-visible service invocation over NDN. In the proposed
UAV application, the drone app can still be viewed as a service container, but
external interactions such as ground-station commands, telemetry discovery,
video retrieval, and cross-drone collaboration use NDNSF service names and
data-centric verification. This distinction is important: NDNSF is not a
replacement for flight controllers or robotics middleware, but a network
service framework that can connect such components across dynamic NDN
environments.

\section{Message Queues and Event Systems}

Message queues and publish-subscribe systems decouple producers and consumers
in time~\cite{messagequeue}. They are useful for telemetry, logs,
asynchronous workflows, and elastic back-end processing. This is relevant to
NDNSF because many service applications include both request-response actions
and event streams. For example, a drone may receive a command as a service
request while publishing telemetry or video metadata asynchronously. A
distributed inference pipeline may publish intermediate results and execution
events while serving user requests.

However, many message queue deployments rely on brokers or logically
centralized infrastructure. Security is often enforced by broker access control
and channel encryption rather than by data objects that remain verifiable
wherever they are stored or forwarded. For mobile or disconnected edge
environments, broker availability and placement can become a limiting
assumption. Moreover, queues do not by themselves define how a user should
choose among multiple providers, verify that a provider is authorized to serve
a specific service, or retrieve large service artifacts as independently
verifiable data.

\section{NDN Computation and Service Frameworks}

NDN already provides useful primitives for service-oriented systems:
hierarchical names, receiver-driven retrieval, in-network caching, adaptive
forwarding, and data signatures~\cite{ndn}. A service framework over NDN must
bridge the gap between these primitives and application-level service
semantics. Prior NDN work explores several regions of this design space.

\subsection{Function-Centric Computation}

Named Function Networking (NFN) treats computation as a first-class named
object by embedding function expressions in names~\cite{nfn}. This model is
powerful because it allows computation and data retrieval to be composed within
the naming system itself. It also suggests a deep integration between content
names and computation names, which is attractive for in-network processing.

The trade-off is complexity. Function-centric naming can make names highly
expressive, but it also makes execution semantics, trust, caching, and
deployment more difficult to reason about. UAV control, live video, and
distributed AI inference require explicit service authorization, provider
selection, long-running state, and application-defined execution policies.
NDNSF therefore chooses a service-invocation abstraction rather than a fully
general function-expression abstraction.

\subsection{Serverless and Edge-Oriented NDN Frameworks}

NFaaS extends NDN toward serverless computing by treating functions as
deployable units that may be executed at edge locations~\cite{nfaas}. Other
edge-oriented frameworks, such as CFEC for information-centric Internet of
Vehicles, explore low-latency microservice execution and dynamic placement in
edge environments~\cite{cfec}. These systems highlight an important direction:
service execution should move closer to users, data sources, or resource-rich
edge nodes.

The open issue is that dynamic placement is only one part of service-oriented
application support. Applications also need a stable API, authorization for
both users and providers, multi-provider selection, large-object handling, and
clear deployment procedures. NDNSF focuses on these runtime service semantics.
It can coexist with future placement mechanisms; for example, a placement
algorithm could decide which drones, storage components, or inference providers
should advertise a given NDNSF service.

\subsection{Remote Invocation Frameworks}

RICE provides remote method invocation in ICN using mechanisms such as
handshakes and thunks to represent computation parameters and
results~\cite{rice}. NSC proposes Named Service Calls for NDN and offers an
RPC-like abstraction for remote execution~\cite{nsc}. These systems are close
to NDNSF's goals because they directly address service invocation rather than
only content retrieval.

The key limitation is that they do not provide the complete set of mechanisms
needed by the application classes targeted in this dissertation. NSC is
primarily one-to-one: a caller contacts an executor, and the protocol is not
designed around simultaneous provider discovery, provider ACK metadata, or
multi-provider selection. RICE provides stronger invocation semantics but
requires changes to the forwarding plane. NDNSF keeps the framework at the
application/runtime layer and uses NDN Data publication and Sync to support
one-to-many service request dissemination without modifying NFD.

\subsection{Application-Specific NDN Frameworks}

Other NDN systems focus on particular domains. DNMP builds a secure network
measurement framework using NDN synchronization and hierarchical
topics~\cite{dnmp}. MIA-NDN explores microservice-centric Interest aggregation
for NDN~\cite{miandn}. SECaaS studies security-as-a-service in multi-access
edge environments~\cite{secaas}. These projects demonstrate that NDN can be
adapted to practical service-like workloads, especially when names and
application semantics are designed together.

The limitation is generality. Application-specific frameworks often make
assumptions about a particular data type, service domain, aggregation pattern,
or deployment topology. NDNSF instead aims to provide a reusable framework
layer that can support different application services, including UAV control,
video, object detection, distributed inference, storage, and collaboration.
The proposed UAV and DistributedInference applications are intended to test
whether the framework is general enough for multiple domains rather than only
one carefully tuned application.

\subsection{Security and Access-Control Mechanisms}

NDN signs Data packets by design, but signatures alone do not answer all
service questions. A service framework must decide who may request a service,
who may provide a service, what data should be encrypted, how keys should be
distributed, and how replayed messages should be rejected. Prior work on
attribute-based encryption and NAC-ABE shows how fine-grained access control
can be integrated with NDN data names~\cite{goyal2006abe,dulal2023nacabe}. These ideas
are central to NDNSF.

The open issue is integration with the service workflow. If access control is
bolted onto individual applications, every application must separately handle
permission retrieval, encryption attributes, message verification, and replay
protection. NDNSF makes these part of the runtime: ServiceController
distributes permissions, service messages are protected according to
service-level attributes, and one-time user/provider tokens bind ACK,
selection, targeted fast-path, and response messages to the current invocation.

\section{NDN-Based Service Systems}

NDN Service Call (NSC)~\cite{nsc} and other NDN application frameworks show
that NDN can support service invocation without IP endpoint binding. NDN's
Interest/Data exchange, name-based routing, in-network caching, and data
signatures address several issues in IP-based systems. These systems are
therefore important evidence that NDN is a promising foundation for services.

The remaining gap is that existing NDN service systems have not provided a
general-purpose runtime that combines:

\begin{itemize}
  \item generic dynamic service APIs,
  \item controller-authorized user and provider permissions,
  \item one-to-many request publication,
  \item provider ACK metadata,
  \item application-defined selection strategies,
  \item adaptive admission control,
  \item large-data references,
  \item provider collaboration, and
  \item service containers that can be evaluated in realistic settings.
\end{itemize}

\section{Design Patterns and Trade-offs}

The related work suggests several recurring design patterns. First, service
frameworks must decouple service names from provider locations. Second, large
inputs and outputs should be represented as named data objects rather than
embedded directly in small request payloads. Third, service state usually
cannot rely only on NDN's Pending Interest Table because service execution may
be long-running, multi-step, or collaborative. Fourth, security must be part of
the invocation workflow rather than an afterthought.

These patterns also expose trade-offs. Function-centric names are expressive
but can be difficult to deploy. Serverless placement is flexible but requires
orchestration. RPC-style invocation is easy to understand but can favor
one-to-one execution. Broker-based systems decouple time but may introduce
infrastructure dependencies. NDNSF's design point is to combine an RPC-like
developer API with NDN's data-centric names, signed Data, Sync-based
publication, and application-visible provider selection.

\section{Open Issues}

Table~\ref{tab:framework-gaps} summarizes the open issues motivating NDNSF.

\begin{table}[H]
\centering
\caption{Open issues in existing service-oriented frameworks.}
\label{tab:framework-gaps}
\begin{tabularx}{\textwidth}{p{0.25\textwidth}X}
\toprule
Issue & Why It Matters \\
\midrule
Endpoint binding & Mobile services and intermittently connected providers
should be invoked by name and capability rather than by stable host address. \\
One-to-many invocation & Reliability and collaboration often require multiple
providers to hear a request and allow the user to select one or more of them. \\
Data-centric security & Data should remain signed, verifiable, and optionally
encrypted at rest, not only while moving through a secure channel. \\
Provider authorization & A user must verify not only that it is allowed to
call a service, but also that the selected provider is authorized to offer it. \\
Provider selection & Provider metadata such as load, queue size, location, and
availability should be part of the service workflow. \\
Large data & Models, video, images, logs, and intermediate tensors should not
be forced into small request payloads. \\
Collaboration & Tasks may require multiple providers, compensation requests,
or aggregation of partial results. \\
Operational validation & Real systems need repeatable certificate, routing,
configuration, and experiment setup to validate the framework outside toy
examples. \\
\bottomrule
\end{tabularx}
\end{table}

NDNSF is designed to address these issues using NDN names, verifiable Data,
sync-based publication, and service-level authorization.

This chapter supports RQ1 by turning the related-work review into a set of
framework requirements rather than a survey checklist. It also prepares RQ2
and RQ3 by identifying which properties must be provided by the NDNSF runtime
instead of being reimplemented separately by each UAV or inference
application.
